\chapter{Introduction}

This thesis-like document, outlines specific contributions, theoretical developments, and experimental investigations on the \textbf{theoretical encoding theory}, logic of operations, \textbf{statistical learning theory} and \textbf{general learning theory}, within the standard outcome also concluded of statistical physics emulations, and explicit specialized movements. Overall, the document posits itself in the \textit{synthesis package}, in which old ideas are reframed, old formalism being look from different angles of perspective, clarification of concepts, restructuralization of different frameworks and their axioms, presumptions, conjectures, hidden dilemma, and structural failure modes, and proposal of newer insights, outlooks, and theories. In specialities, we focus on the \textit{reformation of the theoretical machine-based functions}, of which applies to current practical reduction to computer-dependent choice of construction, \textit{presumption bias} in changing encoding system and the task of encode or decode, hidden assumptions or structural admissions, and ontological framing of different problem sets. On statistical learning, we focus on the proxy of the problem of \textit{double descent} to highlight theoretical-empirical mismatch, theoretical holes and potential failure modes apparent of the old learning theory, reformulation and synthesis of existing frameworks and usages, and novel contribution to a new theory itself, with experimental toy models for investigations. Additionally, we also emphasize the practical notions and selective simulation or emulations encoding on machine systems, plus the general notion of the learning concepts. In conjunction, collaborations and tangent line of thoughts also consider the statistical learning interpretation, of such learning system. 

The core authoring contribution comes from \textbf{Bui Gia Khanh}, from Department of Physics, Hanoi University of Science, Vietnam National University, with experimental results in majority coming from \textbf{Luong Van Tam}, \textbf{Dang Tri Trung} and \textbf{Daud Shahbaz}. Specifically, Luong Van Tam also participated as the main contributor in the chapter about statistical physical interpretation of the learning dynamics. We would like to thank any particular contributors participating in this document and research progress itself. Until further funding is acquired, we would like to clarify of open-source and self-funded research approach on this document. 

\section{Background and motivation}

The new advancements of Artificial Intelligence (AI), and thus techniques in conjunction to it, for example, the \textbf{learning theory}, or the technique of knowledge encoding, the latent space specifications, modelling constructions, agent-theoretic formalism, \textbf{actionable system feedback} techniques, and so on. Within such, resurfaced the utilization of (super)computers and so on, within the baseline of computational theory of construction, and subsequently, computational learning theory / statistical learning theory (CoLT/SLT). Such system relies on the algorithmic approaches and procedural, computable representation of the automata systems to create models with capacities capable of being called as AI, and in a narrower sense, capable of adjusting or self-regulating response upon which it is provided of observations and situation-aware information. Such is where the word \textit{machine learning} comes, and situated itself as the main anecdotal reason or capability of AI, for an artificial intelligence system to be intelligent, is posited to have the learning process. Of a long history in development, ever since the few first systems and classical notions of AI, \cite{mcculloch_pitts_1943,Rosenblatt1958ThePA,minsky_papert_1969,hebb_1949,Robbins:1951,widrow_hoff_1960_adaline,Hestenes:1952} of the old paradigm in symbolic and neuronal approaches of discrete formulaic learning, to the looser \textit{statistical, data-based learning}, particularly as of \cite{10.1145/1968.1972,Vapnik1999-VAPTNO,10.5555/200548,hopfield_1982,fukushima_1980,cybenko_1989,bishop_1995,haykin_1994,kohonen_1982,6797087,nemirovski_yudin_1983,amari_rattray_saad_1998_natural_gradient_online,10.5555/2371238}, and up to modern development of such materials, such as, \cite{Jacot:2018:NTK,khan_bayesian_2024,buschjager_generalized_2020,Manin_2024,devlin_2018_bert,goodfellow2016deep,kaplan2020scalinglawsneurallanguage,Bronstein_2017,bronstein2021geometricdeeplearninggrids,luo1992convergence,tseng2001convergence,friedman2007pathwise,hu_model_2021,wright2015coordinate,janik2021complexitydeepneuralnetworks}, and so on, of the topics about the modern consideration of learning theories, optimization techniques and interpretation landscape, formal languages of the learning processes, as well as architectures on which modern systems are based on; more advanced techniques are developed, for example, as \textit{reinforcement learning}, as seen of \cite{sutton2018reinforcement,bertsekas1996neuro,bertsekas2019dynamic,szepesvari2010algorithms,lattimore2020bandit,kaelbling1996reinforcement,watkins1992q,bellman1957dynamic,williams1992simple,konda2000actor,mnih2015human,mnih2016asynchronous,silver2016mastering,schulman2015trust,schulman2017ppo,lillicrap2015ddpg,haarnoja2018soft,arulkumaran2017brief,ernst2005tree,fu2020d4rl,kostrikov2021offline}, or \textit{online learning} of \cite{shalev2012online,cesa2006prediction,zinkevich2003online,hazan2016introduction} and \cite{littlestone1994weighted} as authoritative examples. Even further, is toward Large Language Model, or LLM (\cite{vaswani2017attention,devlin2019bert,brown2020language,zhao2023llmsurvey,zhao2024llmsurvey2,radford2019language,radford2018improving,raffel2020exploring,touvron2023llama,touvron2023llama2,chowdhery2022palm,ouyang2022training,wei2022chain,kaplan2020scalinglawsneurallanguage,hoffmann2022training,bai2022constitutional}), where cutting edge efforts and so on are to build agentic systems capable of replacing certain labours kind of human in real systems, and so on improving quality of life. The scope is truly large, the field of applications and its limit overwhelmingly extended, and a much richer landscape was formed in those fields.

However, even with such advancements, certain pushback started to emerge. Those are \textit{structural pushbacks} against previously held assumptions, foundational problems regarding concepts taken for granted of hand-waived, or simply is not considered, structural behaviours out of scope of the dynamical trajectory that classical theory tracks (i.e. a theoretical sensory problem), and inadequacy of theory against empiricism, practiced in the modern AI landscapes. Such mismatch highlights substantial loopholes and opaque understanding and patchworks workaround that was present within the field, in which making a large portion of important problems simply too hard to solve analytically or conclusively, as for example, grokking, seen of \cite{power2022grokkinggeneralizationoverfittingsmall}, or \textit{double descent}, as of \cite{belkin_reconciling_2019,nakkiran_deep_2019}, in which exposes the loose nature of the definition on model selections and its landscape, model complexity, and evaluation on model theoretic approaches. The push for information theoretic learning, statistical interpretation, discourses among different interpretation of the learning setting and so is also for artificial intelligence in general, plus the concerns on explainability of artificial intelligence, of which the modern paradigm poorly offer such, as seen of such pushes on \cite{doshi2017towards,lipton2018mythos,ribeiro2016should,lundberg2017shap,molnar2020interpretable}
. Foundational problems still retain its strength unquestioned or unresolved, such as \cite{dreyfus1965alchemy,dreyfus1972what,dreyfus1986mind,suchman1987plans,brooks1991intelligence,searle1980minds,mccarthy1969philosophical,harnad1990symbol}. On the more modern side, of contemporary critiques, \cite{pearl2009causality,marcus2018deep,sutton2019bitter}, and the modernized Chinese Room Arguments. The \textit{frame problem}, while problematic and is often ignored by modern practitioners, often resurfaced when discussions about explainability and operability are in range, as well as the epistemological notions of artificial operatives, or proofs that such artificial constructs truly `understand', or `learn', and its particular computational nature. More than all, is the apparent \textit{theoretical entropy} observed from the landscape of development itself. While empiricism pushes forward, such encouraged particular fields, or rather somewhat all of them, to introduce different analogous interpretation, framing, concepts, or treatments in and of their own to the mechanical details of artificial intelligence. While not solving the inherent foundational problems or assumptive framework that it was based on, various different theories are used to tackle, in a heuristically employed manner, to patch different deviations or unexplainable phenomena together. 

The goal is then to analyse such stack of complexity, to at the foremost, provide a more structural and insightful analysis into the topic of double descent or related phenomena, and the general issue coming around with it. While it might be overstretching, such problem can be considered the ultraviolet catastrophe of learning theory, especially within high-order construct or contrastive claim classical theory seems to hold against newer form of observation. By Kuhnian cycle of theoretical evolution, this signals a shift in theoretical insufficiency, though not catastrophic enough to trigger a theatrical meltdown and total replacement of theory, but rather serious consideration, and dialectic analysis of every moving components in the working machine. Right now, the field is hinging on operational definitions and heuristics, for example that would appear much in the following up chapters, are \textit{model complexity}, and functional semantic of descriptions. Such is to say, a cross-field recognition of \textbf{Rice's theorem}: semantic of operations can not be extracted from a certain given substrate. On more of the technicality, even the definition of \textit{parameterization} is often fiercely debated, such is to observe there exist many modes of definition, in which heterogeneity by variance of structure is considered a feature, but not a categorical danger. Such is one of the stance we will make explicit, which can be right or wrong with the goal to be falsifiable, underlying the report's insights. 

We hope to provide a provisional view onto applying the theory of statistical physics, as seen of widespread usage in the history of development of artificial intelligence and computational method of such, as seen from \cite{hopfield_1982,amit1985spin,gardner1988optimal,seung1992statistical,mezard1987spin,mezard2009information,engel2001statistical,advani2013statistical,bahri2020statmech,mehta2014exact,advani2017high,li2021statmechdeep,cavity_tutorial_delFerraro}, with mentions of such statistical mechanics application in \cite{roberts2021principles}. Some of those statistical learning schemes resulted in the works of G. E. Hinton on the variation of neural network called \textbf{Boltzmann machine}, as seen of \cite{ackley1985learning} and \cite{hinton2012practical}. Nevertheless, such question remains if the statistical scheme overreached of its current applicable form, to indict misleading or only limited improvement and interpretability onto a machine learning and in general an artificial intelligence system dynamics. In addition, assumptions and structural constraints, such as those incurred on Boltzmann machines for it to facilitate an environment interpretable of statistical physics remains a question of interest\footnote{It is widely concern of the fact that there exists many interpretations of machine learning mechanics under different formal language, or application-biased one. Such is, for example, information theoretic, as seen of \cite{RussoZou:2016:InformationTheory,XuRaginsky:2017:InfoGen}, categorical theoretic approaches, of \cite{Manin_2024}, and accordingly, control theory, topological learning, graph theoretical, and statistical physics. While interesting as to illuminate the same problem under different lens, it is troubling of the messiness in such approach, as different theories claim different hypothesis, principles, theorems and axioms. Furthermore, it is of questionable intent to simply forgive the account of applying one's field on another --- there exists the fundamental \textit{domain bias} --- of which as can be seen in statistical physics, for certain event to be applicable of using statistical law, one either must adapt the notion and introduce analogous notion (like temperature and relative model diffusing entropy), or restricting the domain, structure, range of expression, and mechanics --- similar to how Boltzmann machines are. Additionally, applying without care gives increasingly large amount of domain bias, as concepts typically seen and assumed of a system in one theory might interject wrongfully, or simply conflict, narrow down the applicable range of such application, or in the worst case, invalidate the result.}. Even with such problems, we would like to retain the view that statistical physics can be indeed used to interpret double descent, though with such phenomena, requires extreme cautions of logical assumptions and modelling setting. However, running the risk of the field being theoretically fragmented, with each framework claiming explanatory primacy, there is no easy way to fix, but to for now explicitly consider such possibility a reality enforced by observations. 

In all, while not conclusive, the research ultimately aims itself to clarify the epistemological and foundational concepts used of the classical theory. Then, from doing such, we provide a systematic experimental and theoretical clarification of the landscape in which double descent is in, and of the wider theme of classical learning theory friction at large. Doing so, we note that the results are in posture of non-closure, such is to say our works are not \textit{conclusive} or definitive in any restricted form. Hence, sections of the materials are susceptible to reformation, and thus deliberately direct structural and experimental falsification. 
